
@@M0@@-anneaux.
\item
Soit @@M1@@ un anneau, @@M2@@; soit d'autre part @@M3@@ un groupe, opérant trivialement sur @@M4@@~; si on prend pour @@M5@@ le topos des faisceaux d'ensembles sur @@M6@@, sur lesquels @@M7@@ opère, l'anneau @@M8@@ ainsi obtenu est l'anneau des représentations de @@M9@@ dans les @@M10@@-modules projectifs de type fini.
\end{itemize}

2.3. Soit @@M11@@ un anneau commutatif à élément unité~; on se propose de munir le groupe abélien @@M12@@ d'une structure de pré-@@M13@@-anneau, fonctorielle par rapport à @@M14@@. A cet effet, nous allons appliquer les résultats de (1.10) sur le foncteur @@M15@@.

On définit d'abord la multiplication en se donnant un morphisme @@M16@@-bilinéaire de foncteurs~: @@M17@@, par l'élément @@M18@@ de @@M19@@~; cet élément vérifie en effet 1.7.i), et pour voir qu'il vérifie 1.7.ii), il faut montrer que dans @@M20@@, exprimé au moyen des fonctions symétriques élémentaires @@M21@@ des @@M22@@, @@M23@@ ne figure pas dans les coefficients de @@M24@@ pour @@M25@@, ce qui est évident pour des raisons de degré.

La loi ainsi définie sera notée par le signe @@M26@@~; on a donc par définition :
@@M27@@
Elle est distributive par construction, associative et commutative par les diagrammes~:
@@M28@@
